\section{n3. A full conductor-infinite normalization obstruction}
\paragraph{Supports.} G331, n1 and n2.
Work over $k=\mathbb C$ and let
\[
 R=k[x,y,z]_{(x,y,z)}/(xy),\quad S=A_1\times A_2,
 \quad A_1=k[x,z]_{(x,z)},\quad A_2=k[y,z]_{(y,z)},\quad T=k[z]_{(z)}.
\]
These are excellent polynomial localizations; $(xy)=(x)\cap(y)$ proves reduction. Restriction identifies $R=A_1\times_T A_2$: subtract a common line residue from a compatible pair, leaving $xa(x,z)$ and $yb(y,z)$, which lift independently. Denominators with nonzero constant lift to local units. The projections to both factors are onto, so the two idempotents generate $S$ as an $R$-module. Both factors are normal with the component fraction fields. Any integral element in the total fraction ring lies in both normal factors, proving $S$ is the finite normalization. Multiplication by $(1,0)$ forces the common residue of a conductor element to vanish; zero common residue is also sufficient. Hence
\[
 \mathfrak c=(x,y),\quad \mathfrak cS=xA_1\times yA_2,
 \quad R/\mathfrak c=T,\quad S/\mathfrak c=T\times T.
\]
Use $\mathfrak a=\mathfrak m$. The full upstairs boundary has two clopen plane sections. The only conductor-infinite downstairs point is $v_0$ factoring through $T$ as $\operatorname{ord}_z$. Its support $(x,y)$ is NONMINIMAL. Its complete fibre consists of $w_1$ killing $x$ on $A_1$ and $w_2$ killing $y$ on $A_2$, both inducing $\operatorname{ord}_z$. This is the full fibre because a lift killing the conductor factors through the normalized DVR section of $T$. Finite conductor fibres are singletons by $w(s)=v(cs)-v(c)$ for finite-valued conductor $c$.

The continuous obstacle on the ENTIRE section is
\[
 f(v)=-1-1/v(y),\qquad 1/\infty=0.
\]
Here $v(y)\in[1,\infty]$, and reciprocal is continuous on that compact extended interval. Thus $-2\le f\le-1$. It is constant $-1$ on the first plane and $-1-1/w(y)$ on the second, agreeing at the shared fibre.

The regular graded ring of $A_2$ is $k[y,z]$, so $o_2=\operatorname{ord}_{(y,z)}$ is a valuation, normalized, with $f(o_2)=-2$. For EVERY proper primary $J$, writing $d=o_2(J)$ gives $J\subset(y,z)^d$ and hence $w(J)\ge o_2(J)$ for EVERY normalized $w$. Unit ideals satisfy the same inequality. Negation, nonnegative real sums, uniform limits and actual decreasing infima yield
\[
 u(w)\le u(o_2)\quad\hbox{for EVERY }u\in C_{A_2}\text{ or }D_{A_2}. \tag{3}
\]
Thus every second-plane competitor is at most $-2$ everywhere. Constant $-2$ attains that bound; the first-plane envelope is $-1$. The simultaneous choice is the genuine single boundary-power ideal
$J=(x,z)\times(y,z)^2$, containing $(\mathfrak aS)^2$. Its negative profile is $(-1,-2)$, and G331 makes it geometric. Therefore the ENTIRE upstairs $C$ and $D$ envelopes are this admissible continuous tuple. Any base competitor pulls to the second plane, so (3) bounds its value at $v_0$ by $-2$. Constant $-2$ downstairs gives
\[
 P_R^Cf(v_0)=P_R^Df(v_0)=-2,\qquad
 P_S^K(f\rho)(w_1)-P_R^Kf(v_0)=1. \tag{4}
\]
No downstairs global envelope-membership claim is needed. It is not asserted constantly $-2$: $(x,y^2,yz,z^2)$ gives a competitor with first-plane ordinary-order value $-1$ and value $-2$ at $v_0$.

The fibre maximum of $(-1,-2)$ is $-2$ at second-plane monomial points $w_M(y)=M,w_M(z)=1$ but $-1$ at $v_0$. Polynomial lowest-$y$ exponents, with unit denominators, show $w_M\to w_2$ pointwise. Thus that maximum is discontinuous; its second-plane restriction violates (3), so is not even $D$. Direct tensor calculation gives
$S\otimes_R A_1=A_1\times T$ and $S\otimes_R S=A_1\times T\times T\times A_2$. The vertical $T$ components have nonempty boundaries but no field dominating a plane fraction field. They cannot be suppressed in full descent. This is a geometric conductor-infinite obstruction distinct from the curve mixed-zero literal radial obstruction and from the continuous degree-trace obstruction.
